% !TeX program = latexmk --lualatex --interaction=nonstopmode --halt-on-error --output-directory=build --synctex=1

\documentclass{article}
\usepackage[a4paper, margin=30mm]{geometry}

\usepackage[
    colorlinks,
    citecolor=black,
    linkcolor=black,
]{hyperref}
\usepackage[
  style=numeric,
  doi=true, isbn=false, url=false, eprint=false,
  date=year,% print only the year, not the month
]{biblatex}
\addbibresource{references.bib}

\usepackage{amsmath}
\usepackage{amssymb}
\usepackage{cancel}
\usepackage[
    per-mode=fraction
]{siunitx}
\usepackage{tikz}

\newcommand{\dd}{\mathrm{d}}
\newcommand{\DD}{\mathrm{D}}
\newcommand{\pd}[2]{\dfrac{\partial #1}{\partial #2}}
\newcommand{\pdsnd}[2]{\dfrac{\partial^2 #1}{\partial #2^2}}
% \newcommand{\pd}[2]{\partial_{#2} #1}

\setlength{\parindent}{0pt}

\title{Arbitrary Lagrangian--Eulerian Method}
\author{Gidon Bauer}

\begin{document}

\maketitle

\section{Eulerian Formulation of Incompressible Navier--Stokes}

The incompressible Navier--Stokes equations in convective form are
\begin{equation}
    \nabla \cdot \vec{u} = 0
    \label{eq:continuity}
\end{equation}
\begin{equation}
    \pd{\vec{u}}{t} +
    \left(\vec{u} \cdot \nabla\right) \vec{u} +
    \frac{1}{\rho} \nabla p -
    \nu \Delta \vec{u} =
    \vec{f}
    \label{eq:navier-stokes-convective}
\end{equation}
where $\vec{u}$ is the velocity, $\rho$ the mass density, $p$ the pressure, $\nu$ the kinematic viscosity, and $\vec{f}$ an arbitrary momentum source.
This is supplemented by the conservation of an advected and diffused scalar
\begin{equation}
    \pd{s}{t} + \vec{u} \cdot \nabla s - D \Delta s = S
    \label{eq:advection-diffusion-convective}
\end{equation}
with conserved scalar $s$, diffusion coefficient $D$, and source term $S$.

Eq. \eqref{eq:navier-stokes-convective} can be written in conservative form as
\begin{equation}
    \pd{\vec{u}}{t} + \nabla \cdot \left[
        \vec{u} \otimes \vec{u} +
        \frac{1}{\rho} p \mathbb{I} -
        \nu \left(\nabla \vec{u} + \nabla \vec{u}^\top\right)
    \right] = \vec{f},
    \label{eq:navier-stokes-conservative}
\end{equation}
as can Eq. \eqref{eq:advection-diffusion-convective}
\begin{equation}
    \pd{s}{t} + \nabla \cdot \left(s\vec{u} - D \nabla s\right) = S.
    \label{eq:advection-diffusion-conservative}
\end{equation}
The conservative form is preferred for numerical implementation, because the conservation is fulfilled up to machine precision.

\section{ALE Formulation of Incompressible Navier--Stokes}

In the ALE view, we do the calculations on a moving computational grid and not on a stationary one as in the Eulerian view.
The continuity equation Eq. \eqref{eq:continuity} is the same for both views.
The conservation of momentum Eq. \eqref{eq:navier-stokes-convective} changes to \cite{doneaArbitraryAgrangianUlerian2004}
\begin{equation}
    \pd{\vec{u}}{t}\bigg|_\chi +
    \left(\vec{c} \cdot \nabla\right) \vec{u} +
    \frac{1}{\rho} \nabla p -
    \nu \Delta \vec{u} =
    \vec{f}.
    \label{eq:ale-navier-stokes-convective}
\end{equation}
The difference of this equation to Eq. \eqref{eq:navier-stokes-convective} is that the velocity is advected with the relative velocity $\vec{c} = \vec{u} - \vec{w}$, where $\vec{w}$ is the grid velocity and
that the time-derivative is evaluated in the computational domain.
The other operators have to be evaluated in the spatial domain as for the Eulerian view.
This becomes more challenging when the mesh deforms.

The conservation of a scalar Eq. \eqref{eq:advection-diffusion-convective} becomes in the ALE view
\begin{equation}
    \pd{s}{t}\bigg|_\chi + \vec{c} \cdot \nabla s - D \Delta s = S.
    \label{eq:ale-advection-diffusion-convective}
\end{equation}
Also here, the time-derivative is evaluated in the computational domain and the scalar is advected with the relative velocity.

\subsection{Naïve Flux-Based Formulation}
Eq. \eqref{eq:ale-navier-stokes-convective} can also be written in a flux-based form
\begin{equation}
    \pd{\vec{u}}{t}\bigg|_\chi + \nabla \cdot \left[
        \vec{c} \otimes \vec{u} +
        \frac{1}{\rho} p \mathbb{I} -
        \nu \left(\nabla \vec{u} + \nabla \vec{u}^\top\right)
    \right] = \vec{f} - \vec{u} \left(\nabla \cdot \vec{w}\right).
\end{equation}
This introduces a new source term into the equation due to
\begin{equation*}
    \nabla \cdot \left(\vec{c} \otimes \vec{u}\right) =
    \left(\vec{c} \cdot \nabla\right) \vec{u} +
    \vec{u} \left(\nabla \cdot \vec{c}\right),
\end{equation*}
\begin{equation*}
    \nabla \cdot \vec{c} =
    \nabla \cdot \left(\vec{u} - \vec{w}\right) =
    -\nabla \cdot \vec{w},
\end{equation*}
and finally
\begin{equation*}
    \left(\vec{c} \cdot \nabla\right) \vec{u} =
    \nabla \cdot \left(\vec{c} \otimes \vec{u}\right) +
    \vec{u} \left(\nabla \cdot \vec{w}\right).
\end{equation*}
The new source is due to the volume change of the individual cells that occurs when the grid velocity is not divergence-free, i.e. $\nabla \cdot \vec{w} \ne 0$.

The same can be done for Eq. \eqref{eq:ale-advection-diffusion-convective}
\begin{equation}
    \pd{s}{t}\bigg|_\chi + \nabla \cdot \left(
        s \vec{c} -
        D \nabla s
    \right) = S - s \left(\nabla \cdot \vec{w}\right).
\end{equation}

These formulations are similar to the conservative forms Eq. \eqref{eq:navier-stokes-conservative} and \eqref{eq:advection-diffusion-conservative} but they do not conserve discretely due to the new source terms.

\subsection{Conservative Form via the Geometric Conservation Law}

The discrete conservative properties can be recovered by introducing the geometric conservation law \cite{thomasGeometricConservationLaw1979}
\begin{equation}
    \frac{\dd}{\dd t} \int_\Omega J \dd V_\chi = \int_\Omega \left(\nabla \cdot \vec{w}\right) J \dd V_\chi
\end{equation}
with $\dd V = \dd x \dd y \dd z = J \dd \xi \dd \eta \dd \zeta = J \dd V_\chi$ and
$J = \pd{(x, y, z)}{(\xi, \eta, \zeta)}$ being the Jacobian of the transformation.
This can be written in differential form as
\begin{equation}
    \pd{J}{t}\bigg|_\chi - J \nabla \cdot \vec{w} = 0.
    \label{eq:geometric-conservation-law}
\end{equation}

With this, we can modify Eq. \eqref{eq:ale-navier-stokes-convective}
\begin{equation}
    \begin{split}
        \pd{J\vec{u}}{t}\bigg|_\chi =
        J\pd{\vec{u}}{t}\bigg|_\chi + \vec{u}\pd{J}{t}\bigg|_\chi &=
        J \left[
            - \left(\vec{c} \cdot \nabla\right) \vec{u}
            - \frac{1}{\rho} \nabla p
            + \nu \Delta \vec{u}
            + \vec{f}
        \right] +
        J \vec{u} \left(\nabla \cdot \vec{w}\right) \\
        &=
        J \left[
            - \nabla \cdot \left(\vec{c} \otimes \vec{u}\right)
            - \frac{1}{\rho} \nabla p
            + \nu \Delta \vec{u}
            + \vec{f}
        \right] \\
        &=
        -J \nabla \cdot \left[
            \vec{c} \otimes \vec{u}
            + \frac{1}{\rho} p \mathbb{I}
            - \nu \left(\nabla \vec{u} + \nabla \vec{u}^\top\right)
        \right]
        + J \vec{f}.
    \end{split}
\end{equation}

In the discrete case we can then use the following first order update rule
\begin{equation}
    % \begin{split}
    %     &J^{n+1} \vec{u}^{n+1} =
    %     J^{n} \vec{u}^{n} -
    %     \Delta t \left\{
    %         J^{n} \nabla \cdot \left[
    %             \vec{c} \otimes \vec{u}
    %             + \frac{1}{\rho} p \mathbb{I}
    %             - \nu \left(\nabla \vec{u} + \nabla \vec{u}^\top\right)
    %         \right]^{n} -
    %         J^{n} \vec{f}^{n}
    %     \right\} \\
    %     \Rightarrow\ &
        \vec{u}^{n+1} =
        \frac{1}{J^{n+1}} \left(
            J^{n} \vec{u}^{n} -
            \Delta t \left\{
                J^{n} \nabla \cdot \left[
                    \vec{c} \otimes \vec{u}
                    + \frac{1}{\rho} p \mathbb{I}
                    - \nu \left(\nabla \vec{u} + \nabla \vec{u}^\top\right)
                \right]^{n} -
                J^{n} \vec{f}^{n}
            \right\}
        \right)
    % \end{split}
\end{equation}
or alternatively the second order update rule
\begin{equation}
    \begin{split}
        &\vec{u}^{n+\tfrac12} =
        \frac{1}{J^{n+\tfrac12}} \left(
            J^{n} \vec{u}^{n} -
            \frac{\Delta t}{2} \left\{
                J^{n} \nabla \cdot \left[
                    \vec{c} \otimes \vec{u}
                    + \frac{1}{\rho} p \mathbb{I}
                    - \nu \left(\nabla \vec{u} + \nabla \vec{u}^\top\right)
                \right]^{n} -
                J^{n} \vec{f}^{n}
            \right\}
        \right) \\
        &\vec{u}^{n+1} =
        \frac{1}{J^{n+1}} \left(
            J^{n} \vec{u}^{n} -
            \Delta t \left\{
                J^{n+\tfrac12} \nabla \cdot \left[
                    \vec{c} \otimes \vec{u}
                    + \frac{1}{\rho} p \mathbb{I}
                    - \nu \left(\nabla \vec{u} + \nabla \vec{u}^\top\right)
                \right]^{n+\tfrac12} -
                J^{n+\tfrac12} \vec{f}^{n+\tfrac12}
            \right\}
        \right).
    \end{split}
\end{equation}
Note that both require a Poisson equation for the pressure that is not shown here.

The same can be done for Eq. \eqref{eq:ale-advection-diffusion-convective}
\begin{equation}
    \begin{split}
        \pd{Js}{t}\bigg|_\chi =
        J \pd{s}{t}\bigg|_\chi + s \pd{J}{t}\bigg|_\chi &=
        J \left[-\vec{c} \cdot \nabla s + D \Delta s + S\right]
        + J s \left(\nabla \cdot \vec{w}\right) \\
        &= J \left[
            - \nabla \cdot \left(s \vec{c}\right)
            + D \Delta s
            + S
        \right] \\
        &= -J \nabla \cdot \left(s \vec{c} - D \nabla s\right) + JS
    \end{split}
\end{equation}
and we get the same procedure for the first order update rule
\begin{equation}
    s^{n+1} =
    \frac{1}{J^{n+1}} \left\{
        J^n s^n
        - \Delta t \left[
            J^n \nabla \cdot \left(s \vec{c} - D \nabla s\right)^n
            - J^n S^n
        \right]
    \right\}
\end{equation}
and the second order update rule
\begin{equation}
    \begin{split}
        &s^{n+\tfrac12} =
        \frac{1}{J^{n+\tfrac12}} \left\{
            J^n s^n
            - \frac{\Delta t}{2} \left[
                J^n \nabla \cdot \left(s \vec{c} - D \nabla s\right)^n
                - J^n S^n
            \right]
        \right\} \\
        &s^{n+1} =
        \frac{1}{J^{n+1}} \left\{
            J^n s^n
            - \Delta t \left[
                J^{n+\tfrac12} \nabla \cdot \left(s \vec{c} - D \nabla s\right)^{n+\tfrac12}
                - J^{n+\tfrac12} S^{n+\tfrac12}
            \right]
        \right\}.
    \end{split}
\end{equation}

The geometric conservation law has to be discretized in the same way as the other conservation equations for the truncation errors to cancel.

% % = Description of Motion ==========================================================================
% \section{Description of Motion}

% The arbitrary Lagrangian--Eulerian (ALE) method couples Lagrangian methods,
% in which the grid moves with the flow and is deformed, with Eulerian methods,
% in which the grid is constant in time.
% The nodes in the ALE method can be moved with an arbitrary velocity \cite{doneaArbitraryAgrangianUlerian2004}.

% We have three different domains:
% \begin{enumerate}
%     \item The material domain with coordinates $\vec{X}$,
%     \item the spatial domain with coordinates $\vec{x}$, and
%     \item the computational domain with coordinates $\vec{\chi}$.
% \end{enumerate}

% \begin{figure}[ht]
%     \centering
%     \input{figures/material-and-spatial-domain.tikz}
%     \caption{Material and spatial domain and their connection through $\varphi$}
%     \label{fig:material-spatial-domain}
% \end{figure}

% A material point with initial position $\vec{X}$ at $t_0$ will be at point $\vec{x}$ after the transformation $\varphi$.

% \subsection{Lagrangian and Eulerian View}

% In the Lagrangian view, we follow the particles from the initial time:
% \begin{equation}
%     \vec{x}(t) = \varphi(\vec{X}, t).
%     \label{eq:material-to-spatial}
% \end{equation}
% We have the particle velocity and acceleration
% \begin{equation}
%     \vec{v} = \pd{\varphi}{t}(\vec{X}, t) \quad \text{and} \quad
%     \vec{a} = \pdsnd{\varphi}{t}(\vec{X}, t).
% \end{equation}
% The motion is described by the trajectory of the particles \cite{temamMathematicalModeling2005}.
% In the Lagrangian description, there are no convective contributions.
% This means that the material derivative reduces to a time derivative.

% In the Eulerian view, we have the velocity $\vec{u}(x,t)$.
% At each time, $\vec{u}(x,t)$ is the velocity of the particles occupying the position $\vec{x}$.
% The description of the motion is via the velocity field \cite{temamMathematicalModeling2005}.

% \subsubsection{Derivatives}

% We have a particle $\mathcal{P}$ that occupies at time $t_0$ the position $\vec{X}$ and at time $t$ the position $\vec{x}$ is associated with a quantity $f(\mathcal{P}, t)$.
% The function $f$ can be represented in the Lagrangian way by
% \begin{equation}
%     f(\mathcal{P}, t) = g(\vec{X}, t)
% \end{equation}
% or the Eulerian way by
% \begin{equation}
%     f(\mathcal{P}, t) = h(\vec{x}, t).
% \end{equation}
% The two froms are connected via Eq. \eqref{eq:material-to-spatial}.
% \begin{equation}
%     g(\vec{X}, t) = h(\varphi(\vec{X}, t), t)
%     \label{eq:function-lagrangian-eulerian}
% \end{equation}

% We then define the Eulerian derivatives
% \begin{equation}
%     \pd{h}{x_i}(\vec{x}, t) \quad \text{and} \quad \pd{h}{t}(\vec{x}, t)
% \end{equation}
% and the Lagrangian derivatives
% \begin{equation}
%     \pd{g}{X_i}(\vec{X}, t) \quad \text{and} \quad \pd{g}{t}(\vec{X}, t).
% \end{equation}
% The derivative $\pd{g}{t}(\vec{X}, t)$ corresponds to the total derivative $\frac{\DD h}{\DD t}$ \cite{temamMathematicalModeling2005}.

% Via Eq. \eqref{eq:function-lagrangian-eulerian}, we can deduce two properties:
% \begin{enumerate}
%     \item
%     \begin{equation}
%         \pd{g}{X_i}(\vec{X}, t) =
%         % \pd{h(\varphi(\vec{X}, t), t)}{X_i}(\vec{X}, t) =
%         \pd{h(\vec{x}, t)}{x_i}(\vec{x}, t) \cdot \pd{\varphi(\vec{X}, t)}{X_i}(\vec{X}, t)
%     \end{equation}
%     \item
%     \begin{equation}
%         \begin{split}
%             \pd{g}{t} = \frac{\DD h}{\DD t} &= 
%             % \pd{h(\varphi(\vec{X}, t), t)}{t}(\vec{X}, t) =
%             \pd{h(\vec{x}, t)}{t} +
%             \pd{h(\vec{x}, t)}{x_i} \cdot \pd{\varphi(\vec{X}, t)}{t} \\
%             &= \pd{h(\vec{x}, t)}{t} + \pd{h(\vec{x}, t)}{x_i} \cdot u_i
%         \end{split}
%     \end{equation}
% \end{enumerate}

% \subsection{ALE view}

% In the ALE description neither the material configuration nor the spatial configuration are chosen \cite{doneaArbitraryAgrangianUlerian2004}.
% Instead we chose a computational configuration $R_{\vec{\chi}}$.
% This domain moves independently from the material motion and
% can be mapped to both the material and spatial configuration via $\psi$ and $\Phi$.
% \begin{figure}[ht]
%     \centering
%     \begin{tikzpicture}[>=stealth]
%         \begin{scope}[every node/.style={draw, circle, align=center, minimum size=2cm}]
%             \node (material) at (0, 2.5) {Material:\\$R_{\vec{X}}$};
%             \node (spatial)  at (5, 2.5) {Spatial:\\$R_{\vec{x}}$};
%             \node (comp)     at (2.5, 0) {Comp.:\\$R_{\vec{\chi}}$};
%         \end{scope}
%         \draw[->] (material) -- (spatial) node[midway, above] {$\varphi$};
%         \draw[->] (comp) -- (spatial) node[midway, above] {$\Phi$};
%         \draw[->] (comp) -- (material) node[midway, above] {$\psi$};
%     \end{tikzpicture}
%     \caption{Mappings from and to different configurations.}
%     \label{fig:mappings}
% \end{figure}
% In Fig. \ref{fig:mappings} we can clearly see that the particle motion can be expressed as
% \begin{equation}
%     \varphi = \Phi \circ \psi^{-1}
% \end{equation}
% and the mappings are therefore not independent.

% The mapping from the computational to the spatial domain $\Phi$ can be understood as the motion of the grid points in the spatial domain.
% \begin{equation}
%     \begin{gathered}
%         \Phi: R_{\vec{\chi}} \times [0, T) \to R_{\vec{x}} \times [0, T) \\
%         (\vec{\chi}, t) \mapsto \Phi(\vec{\chi}, t) = (\vec{x}, t) \\
%     \end{gathered}
% \end{equation}
% Then the grid velocity is
% \begin{equation}
%     \vec{\hat{v}}(\vec{\chi}, t) = \pd{\vec{x}}{t}\bigg|_{\vec{\chi}}.
% \end{equation}

% The mapping $\psi^{-1}$ transforms the material domain into the computational domain
% \begin{equation}
%     \begin{gathered}
%         \psi^{-1}: R_{\vec{X}} \times [0, T) \to R_{\vec{\chi}} \times [0, T) \\
%         (\vec{X}, t) \mapsto \Phi(\vec{X}, t) = (\vec{\chi}, t) \\
%     \end{gathered}
% \end{equation}
% with corresponding velocity
% \begin{equation}
%     \vec{w} = \pd{\chi}{t}\bigg|_{\vec{X}}
% \end{equation}
% which can be interpreted as the particle velocity in the computational domain.  

% Together this yields
% \begin{equation}
%     \vec{v} = \vec{\hat{v}} + \pd{\vec{x}}{\vec{\chi}} \vec{w}
% \end{equation}
% with the convective velocity
% \begin{equation}
%     \vec{c} := \vec{v} - \vec{\hat{v}} = \pd{\vec{x}}{\vec{\chi}} \vec{w}.
% \end{equation}

% % = Equations in the ALE Formulation ===================================================================
% \section{Equations in the ALE Formulation}

% The equation in the spatial domain
% \begin{equation}
%     \pd{q}{t} + \nabla_{\vec{x}} \cdot \vec{f}(q) = 0
% \end{equation}
% becomes in ALE formulation
% \begin{equation}
%     \pd{J q}{t} + \nabla_{\vec{\chi}} \cdot \left[
%         J A^{-1} \left(\vec{f}(q) - q \vec{w}\right)
%     \right] = 0.
%     \label{eq:general_ale_pde}
% \end{equation}
% The flux $\vec{f}(q)$ must be evaluated in the spatial domain, this has a few implications, for example
% \begin{itemize}
%     \item The derivative in the diffusive flux must account for possible grid rotation.
%     \item Radius appearing in the polar operators must be the correctly move one.
% \end{itemize}
% The vector $\vec{w}$ is the mesh velocity,
% $A = \pd{\vec{x}}{\vec{\chi}}$ is the Jacobian matrix of the deformation and
% $J = \det(A)$ its determinant.
% In two dimensions, factor inside the flux of Eq. \eqref{eq:general_ale_pde} is
% \begin{equation}
%     JA^{-1} =
%     J \left(\begin{matrix}
%         \pd{x}{\xi} & \pd{x}{\eta} \\
%         \pd{y}{\xi} & \pd{y}{\eta} \\
%     \end{matrix}\right)^{-1} =
%     \left(\begin{matrix}
%         \pd{y}{\eta} & -\pd{x}{\eta} \\
%         -\pd{y}{\xi} & \pd{x}{\xi} \\
%     \end{matrix}\right).
% \end{equation}

% The geometric conservation law (GCL)
% \begin{equation}
%     \pd{J}{t} + \nabla_{\vec{\chi}} \left(J A^{-1} \vec{w}\right) = 0
% \end{equation}

% \subsection{Simplifications}

% We assume a constant motion in each spatial direction.
% \begin{equation}
%     \begin{gathered}
%         x(\xi, \eta, t) = x(\xi, \eta, 0) + w_x(t) t \\
%         y(\xi, \eta, t) = y(\xi, \eta, 0) + w_y(t) t
%     \end{gathered}
% \end{equation}
% Initially, the computational grid aligns with the spatial grid.
% \begin{equation}
%     x(\xi, \eta, 0) = \xi \quad \text{and} \quad y(\xi, \eta, 0) = \eta
% \end{equation}

% \begin{equation}
%     \begin{gathered}
%         \pd{}{\xi} x(\xi, \eta, t) =
%         \pd{}{\xi} \left[x(\xi, \eta, 0) + w_x(t) t\right] =
%         \pd{\xi}{\xi} = 
%         1 \\
%         \pd{}{\eta} x(\xi, \eta, t) =
%         \pd{}{\eta} \left[x(\xi, \eta, 0) + w_x(t) t\right] =
%         \pd{\xi}{\eta} =
%         0 \\
%         \pd{}{\xi} y(\xi, \eta, t) =
%         \pd{}{\xi} \left[y(\xi, \eta, 0) + w_y(t) t\right] =
%         \pd{\eta}{\xi} =
%         0 \\
%         \pd{}{\eta} y(\xi, \eta, t) =
%         \pd{}{\eta} \left[y(\xi, \eta, 0) + w_y(t) t\right] =
%         \pd{\eta}{\eta} =
%         1 \\
%         \Rightarrow A = \mathbb{I}
%     \end{gathered}
% \end{equation}

% % = Transformation =================================================================================
% \section{Transformation}

% We define a scalar quantity $f$ in the spatial, computational, and material domain as
% $f(\vec{x}, t)$, $f^{*}(\vec{\chi}, t)$, and $f^{**}(\vec{X}, t)$.

% The spatial description $f(\vec{x}, t)$ can be related to the material description $f^{**}(\vec{X}, t)$ via $\varphi$.
% \begin{equation}
%     f^{**}(\vec{X}, t) = f(\varphi(\vec{X}, t), t) = f \circ \varphi
% \end{equation}
% The gradient then follows as
% \begin{equation}
%     \pd{f^{**}}{(\vec{X}, t)} =
%     \pd{\varphi(\vec{X}, t)}{(\vec{X}, t)} \cdot
%     \pd{f(\vec{x}, t)}{(\vec{x}, t)}
% \end{equation}
% or in matrix form
% \begin{equation}
%     \left(\pd{f^{**}}{\vec{X}} \quad \pd{f^{**}}{t}\right) = 
%     \left(\pd{f}{\vec{x}} \quad \pd{f}{t}\right) \cdot 
%     \left(\begin{matrix}
%         \pd{\vec{x}}{\vec{X}} & \vec{u} \\
%         0 & 1
%     \end{matrix}\right)
% \end{equation}
% From this follows
% \begin{equation}
%     \pd{f^{**}}{\vec{X}} = \pd{f}{\vec{x}} \pd{\vec{x}}{\vec{X}}
% \end{equation}
% and
% \begin{equation}
%     \pd{f^{**}}{t}
%     = \pd{f}{\vec{x}} \cdot \vec{u} + \pd{f}{t}
%     = \pd{f}{t} + (\vec{u} \cdot \nabla) f
% \end{equation}


% ------------------------------------------------------------------------------------------------------

% For the computational we than have
% \begin{equation}
%     f^{**}(\vec{X}, t) = f^{*}((\vec{\chi}, t), t) = f^{*} \circ \psi^{-1}
% \end{equation}
% and with the same argument as above, we get
% \begin{equation}
%     \begin{aligned}
%         &\pd{f^{**}}{t} = \pd{f^*}{t} + \pd{f^*}{t} \cdot \vec{w} \\
%         \Rightarrow &\pd{f^{**}}{t} = \pd{f^*}{t} + \pd{f}{t} \cdot \vec{c} \\
%     \end{aligned}
% \end{equation}

\newpage
\printbibliography

\end{document}